\twocolumn[%
{\begin{center}
\Huge Appendix: Artifact Description/Artifact Evaluation
\end{center}}
]

%%%%%%%%%%%%%%%%%%%%%%%%%%%%%%%%%%%%%%%%%%%%%%%%%%%%%
% Artifact Description (AD)
%%%%%%%%%%%%%%%%%%%%%%%%%%%%%%%%%%%%%%%%%%%%%%%%%%%%%
\appendixAD

\section{Overview of Contributions and Artifacts}

\noindent\textbf{Paper identification.}
This artifact description and evaluation accompanies the PDSW paper
introducing \textsc{PILOT}, a lightweight CPU-based framework for online
prediction and selection of error-bounded lossy-compression configurations for
streaming light-source data. Throughout this form, \textsc{PILOT} is the paper
codename and $A_1$ is its associated computational artifact.

\subsection{Paper's Main Contributions}
The main contributions of \textsc{PILOT} are summarized as follows:
\begin{description}
\item[$C_1$] We characterize compression efficiency and reconstruction quality
across streaming light-source datasets, compressors, relative error bounds,
quality objectives, and consecutive frames.

\item[$C_2$] We introduce lightweight, compressor-specific Ridge models that
use inexpensive frame statistics to predict CR, PSNR, and SSIM and rank
candidate $\langle$compressor, error bound$\rangle$ pairs without exhaustively
compressing every application frame.

\item[$C_3$] We present a lightweight CPU-only online auto-tuning mechanism
that trains on frames from the incoming stream, selects configurations using a
user-defined joint quality--efficiency score, applies the selected configuration
to future frames, and periodically retrains as new frames arrive. Replacing
exhaustive trial-and-error search with model-guided selection achieves up to
$66.64\times$ speedup while exposing a controllable trade-off among tuning
overhead, selection accuracy, and compression quality.
\end{description}

\subsection{Computational Artifact}
\begin{description}
\item[$A_1$] The \textsc{PILOT} artifact contains the model and selection code, five
archived compressor-result CSVs, a compact CSSI ROI figure input, committed
reference run summaries, a manifest
of the 21 paper experiments, a sequential batch runner, result-validation
scripts, five original compression-experiment scripts, an environment
specification, and a notebook that regenerates the
paper's quantitative tables and figures. The development repository is
available at \url{https://github.com/JiaxiangCxx/PILOT_PDSW26}. The immutable
archive DOI will be inserted after artifact deposit.
\end{description}

\begin{center}
\begin{tabular}{rll}
\toprule
Artifact & Contributions & Reproduced evidence \\
\midrule
$A_1$ & $C_1$ & Figures~1--2 \\
      & $C_2$ & Figure~5 and Table~III \\
      & $C_3$ & Figure~6 and Tables~IV--VI \\
\bottomrule
\end{tabular}
\end{center}

%%%%%%%%%%%%%%%%%%%%%%%%%%%%%%%%%%%%%%%%%%%%%%%%%%%%%
\section{Artifact Identification}
%%%%%%%%%%%%%%%%%%%%%%%%%%%%%%%%%%%%%%%%%%%%%%%%%%%%%

\newartifact
\artrel
$A_1$ is the primary evaluation artifact. It supports all three contributions
using archived compression measurements, so evaluators do not need the
original beamline arrays, LibPressio, or the six compressor installations.
The archived measurements under
\texttt{lightsource\_data\_cmp\_results/} are consumed directly; the AD/AE
workflow does not invoke a lossy compressor.
The original compression scripts are included unchanged for reviewer
convenience. Their execution is optional and requires the original datasets,
LibPressio, and the corresponding compressor plugins.

\artexp
The complete workflow executes 21 unique model experiments listed in
\texttt{artifact/paper\_experiment\_manifest.csv}:
\begin{enumerate}
\item five CR--PSNR--SSIM experiments, one for each dataset, with
      $W/P=16/16$ and $\tau=1$;
\item four additional SFC-L objectives (CR, PSNR, SSIM, and CR--PSNR), with
      $W/P=16/16$ and $\tau=1$;
\item seven SFC-L window/granularity experiments: $8/8$, $32/32$, and
      $32/512$ at $\tau=1$, plus $32/512$ at
      $\tau\in\{4,16,64,256\}$; and
\item four additional CR--PSNR experiments on CSSI, FF-HEDM, SFC-GI, and
      XPCS, with $W/P=16/16$ and $\tau=1$.
\item one FF-HEDM PSNR-only experiment with $W/P=16/16$ and $\tau=1$ for
      the rate--quality comparison in Figure~6.
\end{enumerate}
The SFC-L CR--PSNR--SSIM and CR--PSNR $16/16$, $\tau=1$ results are shared
across studies and are each executed only once. Every new result is compared
with its committed reference using frame counts and deterministic
prediction/selection metrics. Runtime is recorded but is not required to match
across hardware.

\arttime
Environment setup takes approximately 5 minutes when the Conda packages are
already cached. On the UNC Charlotte CCI \texttt{cci-abc} reference server, the
complete sequential 21-experiment workflow required exactly
1023.515109~seconds (17.058585~minutes). This measurement includes all model
runs, reference validation, and notebook post-processing, as recorded in
\texttt{artifact/output/ad\_execution\_time.txt}. Runtime may vary with hardware
and load on the shared server.

\artin
\artinpart{Hardware}
The model and analysis require an x86-64 Linux CPU; no GPU is required. The
reference system is the \texttt{cci-abc} server at the College of Computing and
Informatics, University of North Carolina at Charlotte. It contains two Intel
Xeon Silver 4214 processors at 2.20~GHz (12 physical cores per socket, two
threads per core, 24 physical cores and 48 logical CPUs in total) and
251.52~GiB RAM. At least 16~GB RAM and 36~GB free storage are recommended.

\artinpart{Software}
The reference host runs Ubuntu 24.04.4 LTS. The verified project-local
environment uses Python 3.10.20, NumPy 2.2.5, pandas 2.3.3,
scikit-learn 1.7.2, Matplotlib 3.10.6, and SciPy 1.15.3. Exact dependencies are
listed in \texttt{artifact/environment.yml}. No GPU software or compression
library is needed for the evaluated workflow. Git and Git LFS are required to
retrieve the repository and its large SFC-L input CSV.

\artinpart{Datasets / Inputs}
The artifact includes archived CSVs for FF-HEDM
($1440\times2048\times2048$), XPCS ($512\times1813\times1558$), CSSI
($600\times1813\times1558$), SFC-L ($10000\times185\times194$), and SFC-GI
($758\times1440\times1440$). Each row records the frame, compressor, relative
error bound, measured CR, PSNR, SSIM, compression time, and precomputed frame
statistics. These CSVs are the only data inputs used by the batch; compression
is not rerun.

\artinpart{Installation and Deployment}
Clone the repository, retrieve its Git LFS objects, and create the isolated
environment:
\begin{verbatim}
git clone https://github.com/JiaxiangCxx/PILOT_PDSW26.git
cd PILOT_PDSW26
git lfs pull
bash artifact/setup_env.sh
conda activate "$PWD/.conda/pilot-ae"
\end{verbatim}
The runner resolves all Python and Jupyter executables from this activated
project-local environment. No compilation is required.

\artcomp
The workflow is
$T_1\rightarrow T_2\rightarrow T_3\rightarrow T_4\rightarrow T_5$.
$T_1$ validates every archived input and committed reference before starting
the timed run. $T_2$ fits one compressor-specific Ridge model for every active
metric using the current training window. $T_3$ predicts each candidate
$\langle$compressor, error bound$\rangle$ pair, computes its predicted joint
score, and selects the highest-scoring pair. $T_4$ evaluates the selected pair
against the oracle and compares deterministic metrics with the committed
reference. $T_5$ summarizes timings and results and executes
\texttt{model/reproduce\_paper\_results.ipynb} to render and validate the
paper tables and figures.

The CR model uses quantized entropy, SVD truncation fraction, data standard
deviation, relative error bound, and interaction terms. The PSNR model uses
frame range and relative error bound. The SSIM model uses frame range, mean,
variance, and relative error bound and predicts log-transformed SSIM loss.
Unless stated otherwise, the workflow evaluates six compressors, uses
compressor-specific Ridge models with $\alpha=10^{-6}$, equal weights for
active metrics, no-gap windows, and all training frames. Prediction accuracy is
\[
100\,\operatorname{clip}\!\left(
1-\frac{|M_{\mathrm{actual}}-M_{\mathrm{pred}}|}
{|M_{\mathrm{actual}}|},0,1\right).
\]
Selection is correct only when both the selected compressor and relative error
bound match the oracle pair.

\artout
Each run replaces \texttt{artifact/output/datasets/}. Every experiment
directory retains only \texttt{compressor\_timing.csv},
\texttt{joint\_prediction\_accuracy.csv},
\texttt{joint\_selection\_comparison.csv}, and
\texttt{run\_summary.log}. The output root also contains one complete-batch
timing file, \texttt{ad\_execution\_time.txt}. The notebook renders the tables
and figures during execution without saving duplicate derived files.

For SFC-L at $W/P=16/16$ and $\tau=1$, the expected selection accuracies are
89.75\% (CR), 76.17\% (PSNR), 45.29\% (SSIM), 94.59\% (CR--PSNR), and 94.66\%
(CR--PSNR--SSIM). Across five datasets and six compressors, the expected mean
prediction accuracies are 83.35\% for CR, 97.91\% for PSNR, and 76.29\% for
SSIM.

Under the CR--PSNR objective, the expected PILOT aggregate CR values are
341.614, 74.029, 5.044, 4.197, and 10.737 for CSSI, FF-HEDM, SFC-GI, SFC-L, and
XPCS, respectively. For the $32/512$ tuning study, increasing $\tau$ from 1 to
256 reduces the number of decisions from 9,760 to 53 and increases speedup
from 28.59$\times$ to 66.64$\times$, while selection accuracy decreases from
90.73\% to 29.61\% and mean selected actual joint score decreases from 0.677
to 0.623.

%%%%%%%%%%%%%%%%%%%%%%%%%%%%%%%%%%%%%%%%%%%%%%%%%%%%%
% Artifact Evaluation (AE)
%%%%%%%%%%%%%%%%%%%%%%%%%%%%%%%%%%%%%%%%%%%%%%%%%%%%%
\clearpage
\appendixAE

\arteval{1}
\artin
Create the environment as described above. From the artifact root, validate the
archived inputs and list the exact paper experiment matrix:
\begin{verbatim}
python artifact/check_artifact.py
bash artifact/run_all_paper_experiments.sh --list
\end{verbatim}

\artcomp
The first command checks file presence, CSV schemas, row counts, and SHA-256
hashes. The second command reads
\texttt{artifact/paper\_experiment\_manifest.csv} and prints the dataset,
objective, $W/P$, $\tau$, and paper mapping for every model experiment. It does
not execute a model or compressor.

\artout
The validation command must terminate without an exception and print
\texttt{Artifact validation passed}. The list command must report exactly 21
experiments: five dataset-wide CR--PSNR--SSIM runs, four additional SFC-L
objectives, seven SFC-L window/$\tau$ configurations, and four additional
dataset-wide CR--PSNR runs. The SFC-L CR--PSNR run is shared between the
objective study and the aggregate-CR table.

\arteval{2}
\artin
Run the complete model evaluation sequentially so its wall-clock measurement is
interpretable:
\begin{verbatim}
bash artifact/run_all_paper_experiments.sh
\end{verbatim}
The runner resolves \texttt{python} and \texttt{jupyter} from the activated
project-local environment. It replaces the fixed generated-output tree, preventing
stale and new experiment files from being mixed.

\artcomp
The runner reads only the five archived compression CSVs; the notebook also
reads the compact CSSI frame-32 ROI CSV used by Figure~1. For each manifest
entry, it trains the requested compressor-specific Ridge models, evaluates the
no-gap prediction/selection schedule, writes model outputs, and invokes
\texttt{artifact/compare\_reference.py}. It then executes
\texttt{model/reproduce\_paper\_results.ipynb}. Experiments are run
sequentially, and one end-to-end wall-clock duration is recorded for the AD.
No compressor is invoked.

\artout
The final log must state
\texttt{All 21 paper experiments completed and matched their references}.
Confirm that each of the 21 experiment directories contains exactly the four
model files listed in the AD output description. The notebook renders the paper
figures and tables and checks their deterministic values, but does not save
duplicate derived files under \texttt{artifact/output/}. Runtime may vary
across systems; report
\texttt{wall\_time\_s} from \texttt{ad\_execution\_time.txt} as the measured
end-to-end artifact execution time.
